\chapter{Logika kombinacyjna}
\label{ch:komb}

\metryczka{pisać logikę kombinacyjną świadomie: wiedzieć, jaki sprzęt powstaje z
każdej linii, i nie dać się złapać na szerokości wyrażeń ani zatrzaski.}
{2 wieczory}{01-kombinacyjna/}

\section{\texttt{assign} i \texttt{always\_comb}}

Obie konstrukcje opisują logikę kombinacyjną:

\begin{sv}
assign y = sel ? a : b;

always_comb begin
  if (sel) y = a;
  else     y = b;
end
\end{sv}

W \code{always\_comb} używamy przypisania \emph{blokującego} \code{=}. Wewnątrz
bloku działa jak w C, bo kolejne linie widzą wyniki poprzednich, ale cały blok to
\emph{jedna} sieć bramek liczona „natychmiast”. Zmienna pośrednia jest tylko nazwą
dla przewodu:

\begin{sv}
always_comb begin
  t = a ^ b;        // t: nazwa dla przewodu z wyjścia bramki XOR
  y = t & c;
end
\end{sv}

\code{always\_comb} (w odróżnieniu od starego \code{always @(*)}) informuje też
narzędzia o Twoim zamiarze: jeśli wyjdzie z niego coś, co nie jest
kombinacyjne (zatrzask), Yosys odmówi syntezy.

\section{Wektory: wycinanie, sklejanie, powielanie}

\begin{sv}
logic [31:0] x;
x[7:0]                    // najmłodszy bajt
x[31]                     // bit znaku
x[8*i +: 8]               // bajt i: 8 bitów od pozycji 8*i w górę (i może być zmienną)
{a, b}                    // sklejenie: a na starszych bitach
{4{b}}                    // b powtórzone 4 razy
{{20{x[31]}}, x[31:20]}   // rozszerzenie znakiem 12 -> 32 bity (rozdział 4!)
{cout, sum} = a + b;      // sklejenie po lewej: wynik z przeniesieniem
\end{sv}

\section{Liczby ze znakiem w sprzęcie}

Sprzęt nie „wie”, czy wektor bitów jest liczbą ze znakiem. Kod
uzupełnień do dwóch (U2) jest wybrany właśnie dlatego, że \emph{ten sam
sumator} poprawnie dodaje liczby ze znakiem i bez znaku. Znak ma znaczenie
tylko w kilku operacjach:

\begin{center}
\begin{tabular}{@{}lll@{}}
\toprule
Operacja & bez znaku & ze znakiem \\
\midrule
porównanie \code{<} & \code{a < b} & \code{\$signed(a) < \$signed(b)} \\
przesunięcie w prawo & \code{a >> n} (wsuwa 0) & \code{\$signed(a) >>> n} (powiela bit znaku) \\
rozszerzanie do szerszego typu & zerami & kopiami bitu znaku \\
mnożenie (górna połowa), dzielenie & osobne układy & osobne układy \\
\bottomrule
\end{tabular}
\end{center}

Negacja w U2: $-x = \overline{x} + 1$. Stąd odejmowanie $a - b = a +
\overline{b} + 1$, czyli \emph{sumator z negacją jednego wejścia i
przeniesieniem wejściowym 1}. Z tego triku skorzysta ALU w rozdziale 3.

\section{Szerokości wyrażeń: najczęstsze źródło błędów}

W C typ wyrażenia zależy od typów argumentów. W SystemVerilogu \textbf{szerokość
obliczeń zależy też od lewej strony przypisania}:

\begin{sv}
logic [7:0] a, b, s8;
logic [8:0] s9;
s8 = a + b;   // liczone na 8 bitach: przeniesienie ginie
s9 = a + b;   // liczone na 9 bitach: przeniesienie trafia do s9[8]
\end{sv}

\begin{itemize}
  \item Wyrażenie arytmetyczne liczy się w szerokości $\max$(operandy, lewa
    strona). Węższe operandy są rozszerzane.
  \item Liczba bez rozmiaru (\code{5}) ma 32 bity i jest \emph{signed}.
    \code{1'b1} ma 1 bit.
  \item \textbf{Jeśli którykolwiek operand jest unsigned, całe wyrażenie jest
    unsigned.} \code{logic} jest domyślnie unsigned.
\end{itemize}

\begin{pulapka}[\texttt{>>>} bez \texttt{\$signed}]
\code{a >>> n} na wektorze unsigned to zwykłe \code{>>}. Przesunięcie
arytmetyczne wymaga \code{\$signed(a) >>> n}. Podobnie \code{a < b} na
\code{logic} porównuje bez znaku. Ten błąd wychodzi dopiero na liczbach ujemnych,
dlatego testy kursu celowo używają \code{0x80000000} i \code{0xFFFFFFFF}.
\end{pulapka}

Verilator (\code{make lint}) ostrzega o niejawnym obcinaniu i rozszerzaniu
(\code{WIDTHTRUNC}, \code{WIDTHEXPAND}). Traktuj te ostrzeżenia poważnie.

\section{Jaki sprzęt powstaje}

\begin{center}
\begin{tabularx}{\linewidth}{@{}l X l@{}}
\toprule
\textbf{Konstrukcja} & \textbf{Sprzęt} & \textbf{Koszt} \\
\midrule
\code{\& | \^{} \~{}} & bramki, po jednej na bit & 1 poziom \\
\code{c ? a : b}, \code{if/else} & multiplekser 2:1 & tani \\
\code{case} na $N$ wartości & multiplekser $N$:1 (drzewo, $\log_2 N$ poziomów) & umiarkowany \\
\code{a == b} & XOR-y i drzewo OR & tani \\
\code{a + b}, \code{a - b}, \code{a < b} & sumator (łańcuch przeniesień) & średni \\
\code{a << b} ($b$ zmienne) & $\log_2 N$ warstw multiplekserów & spory \\
\code{a * b} & macierz sumatorów & \textbf{drogi} \\
\code{a / b}, \code{a \% b} & sieć odejmowań & \textbf{bardzo drogi} \\
\bottomrule
\end{tabularx}
\end{center}

\paragraph{Ścieżka krytyczna.} Sygnał musi przejść przez najdłuższy łańcuch
bramek w ciągu jednego taktu zegara. Im dłuższy łańcuch, tym niższa
maksymalna częstotliwość. \code{make synth} podaje to jako \emph{longest
topological path}, czyli liczbę bramek na najdłuższej drodze. To przybliżenie, bo
prawdziwe opóźnienia zależą od technologii, ale dobrze pokazuje trendy.

\section{Przykład: sumator ripple-carry}

Sumator pełny (\emph{full adder}) dodaje trzy bity:
\[ s = a \oplus b \oplus c_{in}, \qquad c_{out} = ab + c_{in}(a \oplus b). \]
Łącząc $N$ takich sumatorów łańcuchem przeniesień, dostajemy sumator
$N$-bitowy (rys.~\ref{fig:rca}). Przeniesienie z bitu 0 może wpłynąć na bit
$N-1$, więc ścieżka krytyczna jest proporcjonalna do $N$: $O(n)$.

\begin{figure}[h]
\centering
\begin{tikzpicture}[x=24mm]
  \foreach \i/\n in {0/0, 1/1, 2/2, 4/N\!-\!1} {
    \node[blok, minimum width=14mm, minimum height=10mm] (fa\i) at (-\i,0) {FA};
    \draw[<-, thick] ($(fa\i.north)+(-3mm,0)$) -- ++(0,6mm) node[above left=-1mm, sygnal]{a$_{\n}$};
    \draw[<-, thick] ($(fa\i.north)+(3mm,0)$) -- ++(0,6mm) node[above right=-1mm, sygnal]{b$_{\n}$};
    \draw[->, thick] (fa\i.south) -- ++(0,-6mm) node[below, sygnal]{s$_{\n}$};
  }
  \node at (-3,0) {$\cdots$};
  \draw[<-, thick] (fa0.east) -- ++(6mm,0) node[right, sygnal]{c$_{in}$};
  \draw[->, thick] (fa0.west) -- (fa1.east);
  \draw[->, thick] (fa1.west) -- (fa2.east);
  \draw[->, thick] (fa2.west) -- ($(-3,0)+(5mm,0)$);
  \draw[->, thick] ($(-3,0)+(-5mm,0)$) -- (fa4.east);
  \draw[->, thick] (fa4.west) -- ++(-6mm,0) node[left, sygnal]{c$_{out}$};
  \draw[decorate, decoration={brace, amplitude=5pt, mirror}, akcent] ($(fa4.south west)+(0,-11mm)$) -- node[below=6pt, etykieta]{ścieżka krytyczna: przez wszystkie przeniesienia} ($(fa0.south east)+(0,-11mm)$);
\end{tikzpicture}
\caption{Sumator ripple-carry z $N$ sumatorów pełnych.}
\label{fig:rca}
\end{figure}

Szybsze sumatory (carry-lookahead, Kogge--Stone) liczą przeniesienia
równolegle w $O(\log n)$ poziomach kosztem większej liczby bramek. Na FPGA
operator \code{+} trafia na dedykowane łańcuchy przeniesień, wielokrotnie
szybsze od sumatora złożonego ręcznie z LUT-ów. Wniosek praktyczny:
\textbf{pisz \code{+} i pozwól narzędziom wybrać implementację}. Ręczna
budowa ma sens tylko w celach dydaktycznych albo gdy wiesz, po co.

\section{\texttt{if} kontra \texttt{case} i priorytety}

\begin{sv}
always_comb begin            // łańcuch priorytetów: a wygrywa z b, b z c
  if (a)      y = 1;
  else if (b) y = 2;
  else if (c) y = 3;
  else        y = 0;
end
\end{sv}

Kolejne \code{else if} to kaskada multiplekserów (rys.~\ref{fig:prio}).
Warunek z \emph{pierwszego} \code{if} jest najbliżej wyjścia, więc ma najwyższy
priorytet. Gdy warunki się wykluczają (np. różne wartości jednego sygnału),
\code{case} lepiej oddaje intencję i daje syntezatorowi swobodę zbudowania
płaskiego drzewa.

\begin{figure}[h]
\centering
\begin{tikzpicture}[x=22mm]
  \foreach \i/\c/\v in {1/c/3, 2/b/2, 3/a/1} {
    \pic (m\i) at (\i,0) {muks};
    \draw[<-, thick] (m\i-sel) -- ++(0,-7mm) node[below, sygnal]{\c};
    \draw[<-, thick] (m\i-in1) -- ++(-7mm,0) node[left, sygnal]{\v};
  }
  \draw[<-, thick] (m1-in0) -- ++(-7mm,0) node[left, sygnal]{0};
  \draw[->, thick] (m1-out) -- ++(5mm,0) |- (m2-in0);
  \draw[->, thick] (m2-out) -- ++(5mm,0) |- (m3-in0);
  \draw[->, thick] (m3-out) -- ++(8mm,0) node[right, sygnal]{y};
\end{tikzpicture}
\caption{\code{if / else if / else} jako łańcuch multiplekserów 2:1. Wejście 1 jest wybierane, gdy warunek (nóżka u dołu) jest prawdziwy; warunek \code{a} jest najbliżej wyjścia, więc ma najwyższy priorytet.}
\label{fig:prio}
\end{figure}

\section{Zatrzaski: błąd numer jeden}

\begin{sv}
always_comb begin
  if (en) y = a;    // a gdy en == 0? y ma "pamiętać" starą wartość!
end
\end{sv}

Skoro \code{y} ma pamiętać wartość przy \code{en = 0}, syntezator musi wstawić
element pamiętający, czyli \emph{zatrzask} sterowany poziomem. W układach
synchronicznych to prawie zawsze błąd: psuje analizę czasową, a układ zachowuje
się inaczej na krzemie niż w symulacji.

\textbf{Lekarstwo:} każda zmienna przypisywana w \code{always\_comb} musi dostać
wartość na \emph{każdej} ścieżce. Najprościej przez przypisanie domyślne na
początku bloku:

\begin{sv}
always_comb begin
  y = '0;           // wartość domyślna
  if (en) y = a;
end
\end{sv}

To samo dotyczy \code{case} bez \code{default}, gdy nie wszystkie wartości są
pokryte.

\section{Pętle \texttt{for} i \texttt{generate}: kopiowanie sprzętu}

\begin{sv}
always_comb begin
  parity = 1'b0;
  for (int i = 0; i < 32; i++)
    parity = parity ^ x[i];       // 32 bramki XOR połączone łańcuchem
end
\end{sv}

Pętla jest \textbf{rozwijana w czasie syntezy}: nie ma tu licznika ani iteracji
w czasie, tylko 32 kopie logiki. Liczba iteracji musi być stała.
(Syntezator zwykle przebuduje ten łańcuch w drzewo o głębokości $\log_2 32$, bo XOR
jest łączny.)

Do powielania \emph{instancji} modułów służy \code{generate}:
\begin{sv}
genvar i;
generate
  for (i = 0; i < WIDTH; i++) begin : g_fa      // etykieta bloku jest wymagana
    full_adder fa (.a(a[i]), .b(b[i]), .cin(c[i]), .s(sum[i]), .cout(c[i+1]));
  end
endgenerate
\end{sv}

\section{Przykład: przesuwnik logarytmiczny}

Przesunięcie o zmienną liczbę bitów $0..31$ można zbudować z 5 warstw
multiplekserów: warstwa $k$ przesuwa o $2^k$ albo nie, zależnie od bitu
\code{shamt[k]} (rys.~\ref{fig:barrel}). Każda liczba $0..31$ jest sumą
niektórych potęg dwójki. Przesunięcie w lewo uzyskuje się, odwracając
kolejność bitów, przesuwając w prawo i odwracając z powrotem.

\begin{figure}[h]
\centering
\begin{tikzpicture}[x=19mm]
  \node[sygnal] (in) at (0,0) {a};
  \foreach \k/\s in {1/1, 2/2, 3/4, 4/8, 5/16} {
    \node[blok, minimum width=13mm, minimum height=9mm, font=\sffamily\scriptsize] (l\k) at (\k,0) {$\gg\!\s$\\lub nic};
    \draw[<-, thick] (l\k.south) -- ++(0,-5mm) node[below, sygnal]{shamt[\the\numexpr\k-1\relax]};
  }
  \draw[->, thick] (in) -- (l1);
  \foreach \k [evaluate=\k as \j using int(\k+1)] in {1,...,4} { \draw[->, thick] (l\k) -- (l\j); }
  \draw[->, thick] (l5) -- ++(13mm,0) node[right, sygnal]{y};
\end{tikzpicture}
\caption{Barrel shifter: 5 warstw po 32 multipleksery 2:1.}
\label{fig:barrel}
\end{figure}

\section{Znane ograniczenia Icarus Verilog 12}

Icarus nie obsługuje całego SystemVeriloga. W RTL rzadko to przeszkadza, ale
warto wiedzieć:
\begin{itemize}
  \item brak \code{break} i \code{continue} w pętlach,
  \item brak inicjalizacji całej tablicy w deklaracji (\code{logic [7:0] t [4] = '\{...\}}),
  \item \code{unique case} i \code{priority case} są akceptowane, ale ignorowane,
  \item komunikat \emph{sorry: constant selects in always\_* processes} to tylko
    ostrzeżenie (Makefile go ukrywa).
\end{itemize}

\section{Zadania}

\begin{zadanie}{1.1 --- Multiplekser (\code{rtl/mux4.sv})}
Dwie wersje: \code{always\_comb} z \code{case} oraz jeden \code{assign} z
operatorem \code{?:}. Test sprawdza instancje 8- i 16-bitową, czyli też to, czy
parametr działa.
\end{zadanie}

\begin{zadanie}{1.2 --- Enkoder priorytetowy (\code{rtl/prio\_enc.sv})}
\code{idx} = numer najstarszego ustawionego bitu, \code{valid} = czy jakikolwiek.
Użyj pętli \code{for}. Pętla od bitu 0 do 7, w której każde trafienie nadpisuje
\code{idx}, daje regułę „najstarszy wygrywa”. Dlaczego? Narysuj powstały
łańcuch multiplekserów.
\end{zadanie}

\begin{zadanie}{1.3 --- Sumator ripple-carry (\code{rtl/adder.sv})}
Najpierw \code{full\_adder} z samych bramek, potem \code{adder} przez
\code{generate}, bez operatora \code{+}. Test sprawdza \emph{wszystkie}
kombinacje dla 8 bitów ($2^{17}$ przypadków) i losowe dla 32. Potem
\code{make synth T=adder}: zanotuj liczbę komórek i ścieżkę, zastąp ciało
modułu wersją \code{assign \{cout, sum\} = a + b + cin;} i porównaj.
\end{zadanie}

\begin{zadanie}{1.4 --- Barrel shifter (\code{rtl/shifter.sv})}
SLL, SRL i SRA o \code{shamt} z 5 warstw multiplekserów, bez operatorów
przesunięcia. Ile multiplekserów 2:1 ma Twój układ (bez odwracania)? Porównaj z
\code{make synth T=shifter}.
\end{zadanie}

\begin{zadanie}{1.5 --- Polowanie na zatrzask (\code{rtl/hex7seg.sv})}
Dekoder 7-segmentowy ma dwa błędy prowadzące do zatrzasku. Zanim cokolwiek
poprawisz, uruchom \code{make lint}, \code{make synth T=hex7seg} i
\code{make unit T=hex7seg}. Zauważ, że lint znajduje tylko jeden z dwóch
błędów, a Yosys odmawia syntezy z komunikatem \emph{Latch inferred}. Potem
popraw kod i dokończ tablicę A--F.
\end{zadanie}

\begin{gotowe}
\code{make test} → 5 × PASS, \code{make lint} → wszystko OK.
\end{gotowe}

\begin{kontrolne}
  \item Ile bitów ma wynik \code{a + b} dla 8-bitowych \code{a}, \code{b} przy przypisaniu do 16-bitowej zmiennej? A do 8-bitowej?
  \item Co da \code{\$signed(8'hF0) >>> 2}, a co \code{8'hF0 >>> 2}?
  \item Dlaczego \code{if} bez \code{else} w \code{always\_comb} tworzy zatrzask, a w \code{always\_ff} nie?
  \item Czym różni się pętla \code{for} w RTL od pętli w C? Co jest stałe w czasie syntezy?
  \item Jak głęboka (w poziomach logiki) jest ścieżka krytyczna sumatora ripple-carry dla $N$ bitów, a jak przesuwnika logarytmicznego?
\end{kontrolne}
